%%% SECTION 12: FEDERATED LEARNING FOR PHYSICAL AI ONCOLOGY TRIALS %%%

[This section should be approximately 8-10 pages. The major theme is that
Physical AI Federated Learning \cite{fl-paper} provides a privacy-preserving
framework for unifying Physical AI oncology trials across multiple sites,
enabling coordinated learning and continuous improvement without centralizing
sensitive patient data. This framework addresses a fundamental challenge of
nationwide Physical AI deployment: how to leverage data from all trial sites
for model improvement while maintaining strict patient privacy under HIPAA
\cite{hipaa} and state privacy laws.]

[PRIMARY SOURCE: national-platform/federated\_learning/ directory containing
four chunked files: main\_chunk1\_preamble\_intro\_methods.tex (lines 1-215),
main\_chunk2\_results\_architecture\_pillars\_examples.tex (lines 216-533),
main\_chunk3\_results\_peerreview\_trust\_analytics.tex (lines 534-776), and
main\_chunk4\_discussion\_limitations\_conclusion.tex (lines 777-930). All four
files must be used together to build this section.]

\subsection{The Privacy Challenge in Multi-Site Physical AI Trials}
\label{subsec:fl-privacy}

[Build from national-platform/federated\_learning/main\_chunk1\_preamble\_intro\_methods.tex,
the introduction section. Detail the privacy challenge: Physical AI systems
generate vast amounts of patient data (robot telemetry, sensor readings,
imaging data, treatment outcomes) that could dramatically improve AI models
if pooled, but centralized data collection creates unacceptable privacy risks
and regulatory compliance challenges. Reference federated learning in
healthcare \cite{federated-learning-healthcare} for the broader context of
privacy-preserving machine learning in medical settings.]

\subsection{Development Process and Methodology}
\label{subsec:fl-methods}

[Build from national-platform/federated\_learning/main\_chunk1\_preamble\_intro\_methods.tex,
Section 2 (Methods). Cover the development process, prompt-driven development
methodology, and how the federated learning framework was designed and
validated. Reference Claude Opus 4.6 \cite{claude-opus} and the development
tools used.]

\subsection{Architecture Overview}
\label{subsec:fl-architecture}

[Build from national-platform/federated\_learning/main\_chunk2\_results\_architecture\_pillars\_examples.tex,
the Architecture Overview section. Detail the complete federated learning
architecture including client nodes, aggregation servers, communication
protocols, and security layers. Explain how the architecture maps to the
Physical AI trial site structure from Section~\ref{sec:site-establishment}.]

\subsection{Five Foundational Pillars}
\label{subsec:fl-pillars}

[Build from national-platform/federated\_learning/main\_chunk2\_results\_architecture\_pillars\_examples.tex,
Pillars 1-5. Detail each of the five foundational pillars of the federated
learning framework. For each pillar, explain its role in supporting
Physical AI oncology trials and how it interacts with the other pillars.]

[MAJOR THEME: The five pillars ensure that federated learning in Physical AI
trials is not just technically sound but also regulatory-compliant,
clinically meaningful, privacy-preserving, and operationally practical.
This comprehensive pillar structure differentiates the framework from
generic federated learning implementations.]

\subsection{Workflow Demonstrations}
\label{subsec:fl-workflows}

[Build from national-platform/federated\_learning/main\_chunk2\_results\_architecture\_pillars\_examples.tex,
the Workflow Demonstrations section covering domain, agentic AI, and Physical AI
examples. Detail the three categories of workflow demonstrations showing how
the framework operates in practice. The Physical AI examples are most relevant
to this National Platform and should be given the most emphasis.]

\subsection{Triple AI Peer Review Pipeline}
\label{subsec:fl-peer-review}

[Build from national-platform/federated\_learning/main\_chunk3\_results\_peerreview\_trust\_analytics.tex,
Section 3.7 (Triple AI Peer Review Pipeline). Detail the three-stage AI
peer review process that validates federated learning model updates before
deployment. This mechanism ensures that model improvements are safe and
clinically valid before being applied to Physical AI systems across trial sites.
Reference Claude Opus 4.6 \cite{claude-opus}, GPT-5.4 \cite{gpt-5-4}, and
Gemini \cite{gemini} as the AI systems involved in peer review.]

\subsection{Code Trust and Safety}
\label{subsec:fl-trust}

[Build from national-platform/federated\_learning/main\_chunk3\_results\_peerreview\_trust\_analytics.tex,
the Code Trust and Safety section. Cover how the federated learning framework
ensures code integrity, prevents malicious model updates, and maintains
audit trails for all model changes. Connect to the safety requirements in
Physical AI Adaption: 21 CFR Part 50 \cite{cfr50-adapt} and Physical AI
Adaption: 21 CFR Part 312 \cite{cfr312-adapt}.]

\subsection{Clinical Analytics and Digital Twins}
\label{subsec:fl-analytics}

[Build from national-platform/federated\_learning/main\_chunk3\_results\_peerreview\_trust\_analytics.tex,
the Clinical Analytics and Digital Twins sections. Cover how federated
learning enables cross-site clinical analytics without data centralization,
and how digital twin models are improved through federated updates. Reference
the main repository \cite{main-repo} for digital twin implementations.]

\subsection{Technology Stack and CLI Tools}
\label{subsec:fl-tech-stack}

[Build from national-platform/federated\_learning/main\_chunk3\_results\_peerreview\_trust\_analytics.tex,
the Technology Stack and CLI Tools sections. Detail the software infrastructure
supporting the federated learning framework. Reference the FL repository
\cite{fl-repo} for implementation details.]

\subsection{Integration with National MCP Infrastructure}
\label{subsec:fl-mcp-integration}

[Synthesize how federated learning integrates with the Physical AI National
MCP Servers \cite{national-mcp-paper} described in Section~\ref{sec:national-mcp}.
Cover data flow between MCP servers and federated learning clients, model
update distribution through MCP channels, and coordinated monitoring across
both infrastructure layers. Reference the MCP repository \cite{mcp-repo} for
integration implementations.]

\subsection{Limitations and Future Work}
\label{subsec:fl-limitations}

[Build from national-platform/federated\_learning/main\_chunk4\_discussion\_limitations\_conclusion.tex,
Section 5 (Limitations and Future Work). Cover acknowledged limitations and
planned improvements to the federated learning framework.]
